\newpage
\appendix
\onecolumn
\section{Appendix}

This appendix provides supplementary information to the main paper. The contents are organized as follows:
\begin{itemize}[nosep]
    \item Section~\ref{appendix: hyperparameter setting} details the core hyperparameter setting used in the HELP framework.
    \item Section~\ref{appendix: different backbones} presents additional experimental results across various retrieval baselines using the \textit{Qwen3-30B-A3B-Instruct-2507} backbone.
    \item Section~\ref{appendix: baselines} describes the implementation details and reproduction alignment of the baseline methods to ensure a fair comparison.    \item Section~\ref{appendix: case study} provides a qualitative analysis via a detailed case study on multi-hop reasoning performance.
\end{itemize}

\subsection{HyperParameter Setting}\label{appendix: hyperparameter setting}
The core hyperparameters for HELP were configured as follows.
\begin{table}[h]
    \centering
    \caption{Hyperparameter settings for the HELP framework.}
    \label{tab:hyperparameters}
    \small
    \renewcommand{\arraystretch}{1.2}
    \begin{tabular}{clc}
        \toprule
        \textbf{Symbol} & \textbf{Description} & \textbf{Value} \\
        \midrule
        $N$ & Number of expansion hops & 2 \\
        $n$ & Initial seed triple size & 3 \\
        $k$ & Hypernode beam search size & 50 \\
        $M$ & Quota for logical path-guided passages & 4 \\
        $L$ & Total context size ($M$ + Dense passages) & 5 \\
        \bottomrule
    \end{tabular}
\end{table}

\subsection{Different LLM Backbones}\label{appendix: different backbones}

To further validate that the performance gains of HELP are architectural rather than dependent on a specific language model, we conducted extensive evaluations using \textit{Qwen3-30B-A3B-Instruct-2507} as the generation backbone. The results, summarized in Table~\ref{tab:results_qwen}, demonstrate that HELP consistently outperforms recent graph-based retrieval methods across both simple and complex reasoning tasks.

% Specifically, in \textbf{Simple QA} tasks (NQ and PopQA), \textbf{HELP} maintains a significant lead, suggesting that our HyperNode expansion strategy provides a cleaner and more precise context even for single-hop queries. In the more challenging \textbf{Multi-Hop QA} benchmarks, particularly \textit{2Wiki} and \textit{HotpotQA}, \textbf{HELP} achieves a noticeable improvement in F1 scores compared to \textbf{HippoRAG2} and \textbf{HyperGraphRAG}. This underscores the efficacy of using initial triples as seed HyperNodes to anchor the reasoning path, which effectively mitigates the semantic drift often encountered in traditional linear or simple graph retrieval.

% Notably, on the \textbf{LV-Eval} long-context benchmark, \textbf{HELP} achieves a 44\% relative improvement in EM score over the strongest baseline. This result is particularly meaningful as it indicates that our logical path-guided passage selection (controlled by quota $M$) is more efficient at capturing long-range dependencies than dense retrieval or unguided graph search. The consistent performance across different task types and LLM backbones confirms that \textbf{HELP} provides a robust and generalized framework for logical retrieval-augmented generation.

\begin{table*}[h]
    \centering
    \caption{Performance comparison (EM / F1 scores) of various retrieval methods across different tasks, with all models evaluated using \textit{Qwen3-30B-A3B-Instruct-2507} as the backbone. In each cell, the first value represents the Exact Match (EM) score, and the second represents the F1 score. The best results are in \textbf{bold}. The $^\dagger$ symbol indicates a failure in graph construction due to large corpus size.}
    \label{tab:results_qwen}
    \small
    \vspace{-5pt}
%    \resizebox{\textwidth}{!}{
    \begin{tabular}{l cccccc c}
        \toprule
        & \multicolumn{2}{c}{Simple QA} & \multicolumn{4}{c}{Multi-Hop QA} & \\
        \cmidrule(lr){2-3} \cmidrule(lr){4-7}
        Retrieval & NQ & PopQA & MuSiQue & 2Wiki & HotpotQA & LV-Eval & Avg \\
        \midrule
        \multicolumn{8}{c}{\textit{Graph-based RAG Methods}} \\
        \midrule
        HippoRAG2 \cite{gutierrez2025rag} & 40.4 / 54.4 & 42.2 / 56.1 & \textbf{34.3} / \textbf{45.8} & 60.2 / 69.4 & 58.4 / 72.9 & 7.26 / 9.67 & 40.5 / 51.4 \\
        LinearRAG \cite{zhuang2025linearrag} & 31.6 / 43.8 & 6.6 / 21.9 & 18.7 / 28.1 & 42.4 / 53.2 & 39.6 / 52.4 & 6.5/ 10.7 & 24.2 / 35.0 \\
        HyperGraphRAG \cite{luo2025hypergraphrag} & 30.3 / 42.1 & 7.9 / 22.5 & 22.7 / 36.9 & 45.5 / 58.6 & 54.3 / 70.2 & N/A$^\dagger$ & - \\
        \midrule
        \rowcolor{blue!10} \textbf{HELP (Ours)} & \textbf{43.0} / \textbf{56.9} & \textbf{44.5} / \textbf{57.4} & 33.7 / 44.7 & \textbf{62.0} / \textbf{69.5} & \textbf{60.4} / \textbf{73.9} & \textbf{10.5} / \textbf{12.9} & \textbf{42.4} / \textbf{52.6} \\
        \bottomrule
    \end{tabular}
\end{table*}

\subsection{Baseline Implementation Details}\label{appendix: baselines}
To ensure a fair and comprehensive comparison, all baseline methods we reproduced were evaluated under a strictly controlled environment. Specifically, we aligned these baselines with HELP by using an identical corpus, evaluation datasets, and the same LLM backbones for answer generation. Beyond these necessary alignments, all other internal configurations were kept strictly consistent with their original official implementations as reported in their respective literature. This ensures that the observed performance gains of HELP are attributable to its structural and logical innovations rather than disparate parameter tuning or modified baseline settings.


\newpage
\subsection{Case Study}\label{appendix: case study}

For this case study, HELP is configured with the hyperparameters specified in Table~\ref{tab:hyperparameters} and utilizes \textit{Llama-3.3-70B-Instruct} as the underlying LLM backbone. Table~\ref{tab:case_study_analysis} presents a qualitative comparison of multi-hop reasoning between HELP and several baselines. The query requires a 2-hop link: Princess Elene of Georgia $\to$ Solomon II of Imereti $\to$ Prince Archil of Imereti.

LinearRAG and HyperGraphRAG fail due to semantic distractions, retrieving irrelevant information. While HippoRAG2 retrieves the correct passages, it fails to produce the final answer. In contrast, our method, HELP, anchors the reasoning process by expanding from the Initial Triple Seed HyperNodes. This strategy effectively filters out distractor entities and successfully reconstructs the reasoning path to identify the final answer.

% --- 定义颜色与符号 ---
\definecolor{cgreen}{RGB}{0, 150, 0}
\definecolor{cred}{RGB}{180, 0, 0}
\newcommand{\cmark}{\textcolor{cgreen}{\ding{51}}} % 绿勾
\newcommand{\xmark}{\textcolor{cred}{\ding{55}}}   % 红叉

% --- 定义紧凑列表 ---
\newlist{tabitemize}{itemize}{1}
\setlist[tabitemize]{label=\textbullet, nosep, leftmargin=1em, after=\vspace{-0.5\baselineskip}, before=\vspace{-0.5\baselineskip}}


\begin{table*}[htbp]
    \centering
    \small
    \renewcommand{\arraystretch}{1.3}
    \caption{A case study illustrating how different methods handle a 2-hop query requiring cross-passage inference. HELP successfully leverages seed HyperNodes to anchor the correct reasoning path, whereas other baselines suffer from semantic drift or hallucination.}
    \begin{tabularx}{\textwidth}{
        >{\bfseries}p{0.18\textwidth} % 第一列：标签列，加粗，固定宽度
        >{\RaggedRight\arraybackslash}X % 第二列：内容列，左对齐，自适应宽度
    }
        \toprule
        % === 题目部分 ===
        Question & Who is the husband of Princess Elene Of Georgia? \\
        \midrule
        Ground Truth & Prince Archil of Imereti \\
        \midrule
        Reasoning Chain & $\text{Princess Elene of Georgia} \xrightarrow{\text{mother of}} \text{Solomon II of Imereti} \xrightarrow{\text{had father}} \text{Prince Archil of Imereti}$ \\
        \bottomrule
        \addlinespace[0.5em]

        % === 模型对比部分 ===
        
        LinearRAG & 
        \textbf{Retrieved Passages: } \newline
        1) \xmark \ Grand Duchess Elena Vladimirovna of Russia: ...husband was Prince Nicholas... \newline
        2) \xmark \ Prince Adarnase of Kartli: ...natural son of Levan of Kartli by a concubine... \newline
        3) \xmark \ Princess Charlotte of Württemberg: ...the wife of Grand Duke Michael Pavlovich of Russia... \newline
        4) \xmark \ Clous van Mechelen: ...is a Dutch musician, arranger, and actor... \newline
        5) \xmark \ Bernhard III, Prince of Anhalt-Bernburg: ...was the eldest son of bernhard ii... \newline
        \textbf{Prediction:} \xmark \ There is no direct information provided about the husband of Princess Elene of Georgia in the given passages. \newline \\
        \midrule

        HyperGraphRAG & 
        \textbf{Retrieved Passages: } \newline
        1) \xmark \ Gundobad: ...He was the husband of Caretene. \newline
        2) \xmark \ Eunoë (wife of Bogudes): Eunoë was the wife of Bogudes, King of Mauretania. \newline
        3) \xmark \ Engelbert III of the Mark: ...Adolph was the eldest son of Count Adolph II... \newline
        4) \xmark \ Megingoz of Guelders: ...He married Gerberga of Lorraine... \newline
        5) \xmark \ Princess Rodam of Kartli: ...Princess Rodam married King George VII... \newline
        \textbf{Prediction:} \xmark \ King George VII of Imereti \newline \\
        \midrule

        HippoRAG2 & 
        \textbf{Retrieved Passages: } \newline
        1) \cmark \ Princess Elene of Georgia: ...She was the mother of Solomon II of Imereti... \newline
        2) \xmark \ Grand Duchess Elena Vladimirovna of Russia: ...husband was Prince Nicholas... \newline
        3) \xmark \ Princess Rodam of Kartli: ...Princess Rodam married King George VII... \newline
        4) \cmark \ Solomon II of Imereti: ...He was born as David to Prince Archil of Imereti... \newline
        5) \xmark \ Grand Duchess Elena Pavlovna of Russia: ...his second wife Sophie Dorothea of Württemberg... \newline
        \textbf{Prediction:} \xmark \ None mentioned. \newline \\
        \midrule
        
        HELP & 
        \textbf{Expansion from Initial Triples as Seed HyperNodes} \newline
        ("princess elene of georgia", "daughter of", "heraclius ii") \newline 
        \textcolor{mygreen}{\textbf{("princess elene of georgia", "mother of", "solomon ii of imereti")}} \newline
        ("princess elene of georgia", "born in", "georgia") \newline 
        
        \textbf{Retrieved Passages: } \newline
        1) \cmark \  Princess Elene of Georgia: ...She was the mother of Solomon II of Imereti... \newline
        2) \cmark \  Solomon II of Imereti: ...He was born as David to Prince Archil of Imereti... \newline
        3) \xmark \  Grand Duchess Elena Vladimirovna of Russia: ...husband was Prince Nicholas... \newline
        4) \xmark \  Princess Rodam of Kartli: ...Princess Rodam married King George VII... \newline
        5) \xmark \  Grand Duchess Elena Pavlovna of Russia: ...his second wife Sophie Dorothea of Württemberg... \newline
        
        \textbf{Prediction:} \cmark \ Prince Archil of Imereti \newline \\
        \bottomrule
    \end{tabularx}
    
    \label{tab:case_study_analysis}
\end{table*}